\section{Week 5: Pretraining, GPT-Style Models, and In-Context Learning}

\subsection{Learning Objectives}

By the end of this week, you should be able to:
\begin{itemize}
  \item Explain the GPT-style pretraining paradigm
  \item Understand in-context learning as inference-time behavior
  \item Distinguish memorization from generalization
  \item Identify evaluation pitfalls in language modeling
  \item Reason about why prompting works at all
\end{itemize}

\subsection{LLM Components Covered}

\textbf{Training paradigm}
\begin{itemize}
  \item Decoder-only Transformers
  \item Autoregressive pretraining
  \item Next-token prediction at scale
\end{itemize}

\textbf{Inference behavior}
\begin{itemize}
  \item Prompting
  \item Few-shot and in-context learning
  \item Sampling strategies
\end{itemize}

\subsection{Conceptual Core: Pretraining as Representation Learning}

GPT-style models are trained on massive corpora using a single objective:

\[
\mathcal{L} = -\mathbb{E} \left[ \log P_\theta(x_t \mid x_{<t}) \right]
\]

No task labels are provided.
All downstream behavior emerges from statistical regularities in text.

This leads to a critical conceptual shift:
\begin{quote}
Training produces a general-purpose sequence model; prompting selects behavior.
\end{quote}

\subsection{Decoder-Only Architecture}

GPT-style models use:
\begin{itemize}
  \item causal self-attention,
  \item no encoder,
  \item no explicit task heads.
\end{itemize}

This simplicity enables:
\begin{itemize}
  \item scalable pretraining,
  \item flexible prompting,
  \item uniform treatment of all tasks as text completion.
\end{itemize}

\subsection{In-Context Learning}

In-context learning refers to the ability of a model to adapt its behavior based on examples provided in the prompt, without gradient updates.

Key properties:
\begin{itemize}
  \item occurs at inference time,
  \item does not modify parameters,
  \item is sensitive to prompt format and ordering.
\end{itemize}

The dominant hypothesis is that in-context learning is a form of implicit meta-learning learned during pretraining.

\subsection{Memorization vs Generalization}

Large models inevitably memorize parts of their training data.
However, they also:
\begin{itemize}
  \item recombine patterns,
  \item interpolate between examples,
  \item generalize to novel inputs.
\end{itemize}

Distinguishing these regimes empirically is difficult and error-prone.

\subsection{Evaluation Pitfalls}

Common evaluation failures include:
\begin{itemize}
  \item data contamination,
  \item overly narrow benchmarks,
  \item prompt overfitting,
  \item interpreting fluency as correctness.
\end{itemize}

Careless evaluation can dramatically overestimate capability.

\subsection{Systems Engineering Implications}

\textbf{Requirements}
\begin{itemize}
  \item Clear separation between training and evaluation data
  \item Reproducible prompting protocols
\end{itemize}

\textbf{Architecture}
\begin{itemize}
  \item Prompt templates as system artifacts
  \item Versioning of prompts and decoding parameters
\end{itemize}

\textbf{Verification \& Validation}
\begin{itemize}
  \item Hold-out evaluation
  \item Prompt robustness testing
\end{itemize}

\textbf{Operations}
\begin{itemize}
  \item Monitor prompt drift
  \item Guard against prompt injection
\end{itemize}

\subsection{Human Factors and Failure Modes}

Humans naturally anthropomorphize in-context learning as “understanding.”

Failure modes include:
\begin{itemize}
  \item assuming persistent learning,
  \item overtrusting few-shot performance,
  \item neglecting distribution shift.
\end{itemize}

\subsection{Key References}

\begin{itemize}
  \item \citet{radford2018improving} — GPT
  \item \citet{brown2020language} — GPT-3 and in-context learning
\end{itemize}

\subsection{Recommended University Lectures}

\begin{itemize}
  \item Stanford CS224N — Language models and pretraining  
  \url{https://web.stanford.edu/class/cs224n/}

  \item Stanford CS324 — Foundation models  
  \url{https://stanford-cs324.github.io/winter2022/}
\end{itemize}

\subsection{Practitioner Lab}

See \texttt{labs/week05/week05\_lab\_pretraining\_incontext.ipynb}.

\subsection{Week 5 Takeaway}

\begin{quote}
Pretraining creates a powerful statistical engine; prompting steers it, but does not retrain it.
\end{quote}


